\section[Desenvolvimento do Projeto]{Desenvolvimento do Projeto}

\subsection{Repositório do Código Fonte do Projeto}

Para organizar o desenvolvimento do Plan Line, a equipe optou por separar o projeto em quatro repositórios principais: Wiki, backend, mobile e web. Essa divisão ajudou a manter o projeto mais encapsulado, facilitando a organização das responsabilidades e a consulta das informações técnicas por frente de desenvolvimento.

{\footnotesize
\begin{itemize}
    \item Wiki: \url{https://tools.ages.pucrs.br/planline/wiki}
    \item Backend: \url{https://tools.ages.pucrs.br/planline/backend}
    \item Mobile: \url{https://tools.ages.pucrs.br/planline/mobile}
    \item Web: \url{https://tools.ages.pucrs.br/planline/web}
\end{itemize}
}

A Wiki foi utilizada como fonte de documentação do projeto, concentrando informações sobre arquitetura, banco de dados, protótipos e decisões técnicas. Os demais repositórios separavam as camadas de implementação da solução, o que contribuiu para uma organização mais clara do trabalho.

\subsection{Banco de Dados Utilizado}

O projeto utilizou Firebase para gerenciamento das sessões e PostgreSQL para armazenamento dos dados gerais da aplicação. Essa composição buscava atender às necessidades específicas do sistema, que envolvia sessões de votação em tempo real e persistência estruturada de informações relacionadas às histórias, usuários e resultados.

As informações sobre o banco de dados foram documentadas na Wiki do projeto:

{\footnotesize
\begin{itemize}
    \item Banco de dados: \url{https://tools.ages.pucrs.br/planline/wiki/-/wikis/banco_dados}
\end{itemize}
}

\subsection{Arquitetura Utilizada}

A arquitetura escolhida para o Plan Line foi baseada em Clean Architecture. Essa escolha tinha como objetivo organizar melhor as responsabilidades do sistema, separando regras de negócio, infraestrutura, interfaces e camadas externas. Para uma equipe em formação, essa abordagem também funcionou como uma oportunidade de aprendizado sobre organização de código e divisão de responsabilidades em aplicações web.

As informações de arquitetura foram registradas na Wiki:

{\footnotesize
\begin{itemize}
    \item Arquitetura: \url{https://tools.ages.pucrs.br/planline/wiki/-/wikis/arquitetura}
\end{itemize}
}

\subsection{Protótipos das Telas Desenvolvidas}

Os protótipos das telas foram desenvolvidos no Figma com apoio da designer da Triider. Essa etapa foi importante para alinhar expectativas com os stakeholders e orientar a implementação das interfaces web e mobile. Os mockups também ajudaram a equipe a compreender melhor os fluxos esperados para criação de sessões, participação nas votações e visualização dos resultados.

{\footnotesize
\begin{itemize}
    \item Mockups: \url{https://tools.ages.pucrs.br/planline/wiki/-/wikis/mockups}
\end{itemize}
}

\subsection{Tecnologias Utilizadas}

Entre as principais tecnologias utilizadas no projeto estavam HTML, CSS, JavaScript, TypeScript, React, Tailwind CSS, Node.js, PostgreSQL e Firebase. A equipe também trabalhou com conceitos de desenvolvimento web, integração entre frontend e backend, importação e exportação de arquivos \texttt{.csv}, criação de endpoints e organização de componentes reutilizáveis.

No meu caso, o contato com essas tecnologias foi especialmente relevante porque muitas delas estavam além do meu conhecimento inicial. Durante o projeto, precisei estudar os fundamentos de frontend, compreender a estrutura de aplicações React, aprender a utilizar Tailwind CSS e, posteriormente, iniciar minha aproximação com backend, Node.js, Prisma, PostgreSQL e métodos HTTP.
